% !TeX root = ../复分析笔记.tex
% 由主文档通过 \input 载入；请编译主文档。
\setcounter{chapter}{0} % 保留原书章号。
\chapter{先修工具与常用结论索引}
\label{chap:ca-prerequisites}

本附录列出正文调用的实分析工具，并标明作为外部结果引用的平面拓扑结论。圆盘上的柯西理论通过三角形细分、局部原函数与积分估计证明；推广到一般简单闭曲线时，还需要关于内部、外部与绕数的拓扑事实。将这些依赖明确列出，可以区分局部分析证明与大域拓扑条件。

\section{实分析与紧性工具}

正文使用欧氏空间中的紧性判别、紧集上连续函数的一致连续性和极值存在性，以及连续偏导数保证实可微的结论。全纯反函数定理调用实 $C^1$ 反函数定理，其使用条件是所考察点的雅可比行列式非零；在复可微映射处，这个行列式等于 $|f'|^2$。紧度量空间上的等度连续抽取原理则已在引理~\ref{lem:ca-ascoli} 中给出证明，随后用于蒙泰尔定理。

有限曲线上的一致收敛允许交换积分与极限，其误差由曲线长度乘以一致误差控制；无穷区间上的积分还需尾部控制。伽马函数一章使用的充分条件是存在与参数无关的可积上界：在 $0<t<1$ 与 $t>1$ 分别估计，再应用支配收敛。双重积分换序使用绝对可积性。参数微分可以先在有限区间进行，再验证导数积分在参数紧集上具有统一可积控制，不能仅凭逐点收敛交换这些运算。

\section{简单闭曲线的拓扑工具}
\label{sec:ca-jordan-tool}

开区域的折线连通性见命题~\ref{prop:ca-polygonal-connected}。以下若当曲线定理（Jordan curve theorem）及其可收缩性结论作为外部拓扑工具引用。

\begin{theorem}[简单闭曲线的内外区域]
\label{thm:ca-jordan-curve}
非退化简单闭曲线的补集恰有两个连通分支：一个有界，称为内部；另一个无界，称为外部。两者的边界均为该曲线。该曲线能在任意包含曲线及其内部的开集中收缩为一点。若曲线还分段 $C^1$，则其绕数在外部为零，在内部恒为 $1$ 或恒为 $-1$；内部绕数为 $1$ 的方向称为正向。
\end{theorem}

这一形式可核对 T.~Tao 的\href{https://terrytao.wordpress.com/2016/10/02/math-246a-notes-3-cauchys-theorem-and-its-consequences/}{Math 246A, Notes 3：Cauchy's theorem and its consequences（2016），Theorem 50 及附录}。对自交闭曲线，不能直接使用唯一“内部”的说法，应按各补集分支的绕数计算。对圆周，绕数及收缩已有显式参数化，因此局部圆盘理论不需要借助上述一般结论。

第~\ref{chap:ca-normal} 章提及的卡拉泰奥多里边界定理也作为引用：若 $D$ 是有界若当区域，圆盘与 $D$ 之间的双全纯映射可延拓为闭包之间的同胚。可参见 T.~Tao 的\href{https://terrytao.wordpress.com/2018/04/12/246c-notes-2-circle-packings-conformal-maps-and-quasiconformal-maps/}{Math 246C, Notes 2（2018），Theorem 20}。这一额外的边界结论不用于黎曼映射定理的存在性证明；“内部存在双全纯映射”与“闭包之间存在相应同胚”必须分别检查条件。

\section{关键定理的条件索引}

柯西定理的单连通形式见定理~\ref{thm:ca-cauchy-domain}。区域内的全纯性与闭路可收缩性缺一不可；在去心平面上，$1/z$ 全纯却沿单位圆积分非零。柯西积分公式见定理~\ref{thm:ca-cauchy-formula}，其简单闭曲线版本要求函数在曲线连同内部的开邻域全纯；只在曲线附近全纯不够。

恒等定理见定理~\ref{thm:ca-identity}，零点的聚点必须位于区域内部，区域还须连通。最大模原理见定理~\ref{thm:ca-maximum}，结论针对全纯函数的模；连续函数和一般实可微函数没有同样性质。调和函数的实值最大、最小值原理见定理~\ref{thm:ca-harmonic-max}，不可把实部的极值与复函数模的极值混用。

留数定理见定理~\ref{thm:ca-residue}，必须列出所围区域内全部孤立奇点，并检查积分路径不经过奇点。鲁歇定理见定理~\ref{thm:ca-rouche}，边界上的严格模不等式确保整个扰动过程没有边界零点，结论中的零点数按重数计算。

蒙泰尔定理见定理~\ref{thm:ca-montel}，局部一致有界中的界必须对整个函数族统一。黎曼映射定理见定理~\ref{thm:ca-riemann-mapping}，其区域必须非空、单连通且为复平面的真子域；导数归一化只负责唯一性，不能代替这些存在性条件。单值化定理见定理~\ref{thm:ca-monodromy}，须先有沿每条路径延拓的能力，再用同伦消除路径依赖。

魏尔斯特拉斯与米塔--莱夫勒定理分别见定理~\ref{thm:ca-weierstrass-product}、定理~\ref{thm:ca-mittag-leffler}，指定点集在复平面中不得有有限聚点。龙格定理见定理~\ref{thm:ca-runge}，函数须在紧集的开邻域全纯，允许极点的位置要覆盖补集各有界分支。几个定理都具有存在性结论，但各自控制的数据不同，不能相互直接替代。

\section{证明路线与拓展阅读}

本书的证明主线为：复可微性与三角形细分导出柯西理论；柯西公式导出局部幂级数、刚性与导数估计；洛朗展开和留数导出零点计数；鲁歇定理与紧性导出赫尔维茨和黎曼映射；局部对数与一致收敛导出乘积及主部构造；柯西积分离散化导出龙格逼近。每条路线中的条件均在相应定理处陈述，进阶定理不反过来充当其基础步骤的证明。

进阶部分可配合 MIT \href{https://ocw.mit.edu/courses/18-112-functions-of-a-complex-variable-fall-2008/pages/lecture-notes/}{18.112 官方讲义}阅读：Lecture 17 涉及主部构造，Lecture 18 涉及无穷乘积，Lecture 19 涉及正规族，Lecture 20 讨论黎曼映射证明中的一个易错环节。关于黎曼曲面、几何与其他学科的联系，可参照 C.~McMullen 的\href{https://abel.math.harvard.edu/~ctm/home/text/class/harvard/213a/06/html/home/course/course.pdf}{Harvard Math 213a：Complex Analysis Notes}。这些材料用于核对结论与扩展阅读；正文中的证明保持独立组织，并明确区分本书已证结果和仅作阅读方向的高级专题。

